\section{How is a leak detected, and can the labels be trusted?}
\label{app:labels}

A rule on the tool log finds the first call that writes the secret out, by exact string match, so the main label needs no judge.
An automated judge, used only to describe the leaks, agrees with the author's blind labels (\cref{tab:labels}).

\paragraph{The leak label.}
A call leaks if the protected value appears, in full, in its reply or in the arguments of any of its tool calls.
This covers channel messages, the end-of-turn reply that the harness posts to the channel, and file writes, wherever they write.
The value appearing in private reasoning does not count.
The label is the one used for gpt-oss-20b's frozen confirmation, applied unchanged to Qwen3-32B.
A few Qwen3-32B calls write the value back into the protected file itself.
Excluding them changes the per-call result very little (\cref{app:robustness}).

\paragraph{The label in the Qwen3-32B round-level test.}
The plan for the Qwen3-32B fresh test episodes fixed a slightly different leak.
There, a leak is the complete value posted to the team channel, or written into another file in the five scenarios other than the main one.
Scored calls stop before the secret is first written out anywhere.
We report that test as it was planned.
On gpt-oss-20b, the round-level analyses in \cref{app:gptoss-round} use a channel post as the leak, again with scored calls stopping before the secret is first written out.

\paragraph{How leaks go out.}
On Qwen3-32B, about half of the leaks go out as the ordinary end-of-turn reply and half inside tool arguments.
On gpt-oss-20b, almost every first leak is a file write, typically the whole configuration block copied into the handoff file.
An end-of-turn reply that carries a leak is almost always addressed to the peer, judged by a decision model (\cref{tab:labels}).

\paragraph{Deliberate or accidental.}
A second decision model, Luna Decisions \citep{openai2026luna}, reads the first message that publishes the value.
It labels the leak deliberate if the message offers the real value for use, and accidental if it exposes the value inside a refusal or an illustration.
Most Qwen3-32B leaks are deliberate, and on gpt-oss-20b nearly all are.
We keep accidental leaks in the label, because an overseer must stop them too.
On the eventual label, refitting the probe on deliberate leaks only gives a similar ranking (\cref{app:robustness}).

\paragraph{Peer claims of authority.}
A decision-model pass flagged peer messages that claim to speak for the operator, or to be an authorised recipient.
Leaks were no more common in episodes with such claims (\cref{tab:labels}).

\begin{table}[h]
  \centering
  \caption{Label checks. Agreement compares the judge with the author's blind labels on a sample of leaks.}
  \label{tab:labels}
  \small
  % Generated by generate_tables.py; edit registry table configuration.
\begin{tabular}{lllll}
\toprule
Target & Audit cohort & Check & Yes & Denominator \\
\midrule
Qwen & Historical first channel leak & Jev addressed to peer & \labelAddressOne{} & \labelAddressNOne{} \\
gpt-oss & Historical first channel leak & Jev addressed to peer & \labelAddressFive{} & \labelAddressNFive{} \\
Qwen & Historical episode outcome & Operator claim among leaks & \labelAuthorityQwenOperator{} & \labelAuthorityNQwenOperator{} \\
Qwen & Historical episode outcome & Receiver claim among leaks & \labelAuthorityQwenReceiver{} & \labelAuthorityNQwenReceiver{} \\
GPT & Historical episode outcome & Operator claim among leaks & \labelAuthorityGPTOperator{} & \labelAuthorityNGPTOperator{} \\
GPT & Historical episode outcome & Receiver claim among leaks & \labelAuthorityGPTReceiver{} & \labelAuthorityNGPTReceiver{} \\
\bottomrule
\end{tabular}

\end{table}
